\IfFileExists{stacks-project.cls}{%
\documentclass{stacks-project}
}{%
\documentclass{amsart}
}
\usepackage{amssymb}

% For dealing with references we use the comment environment
\usepackage{verbatim}
\usepackage{fancyvrb}
\usepackage{pmboxdraw}
\newenvironment{reference}{\comment}{\endcomment}
%\newenvironment{reference}{}{}
\newenvironment{slogan}{\comment}{\endcomment}
\newenvironment{history}{\comment}{\endcomment}

% For source figures retained as assets
\usepackage{graphicx}

% For commutative diagrams we use Xy-pic
\usepackage[all]{xy}

% We use 2cell for 2-commutative diagrams.
\xyoption{2cell}
\UseAllTwocells

% We use multicol for the list of chapters between chapters
\usepackage{multicol}

% This is generally recommended for better output
\usepackage{lmodern}
\usepackage[T1]{fontenc}
\usepackage{textalpha}

% For cross-file-references
% We prefer the snapshot tmp/xref/#2.aux,
% which latexmkrc replaces atomically, so
% that documents compiled simultaneously
% never read a half-written .aux file.
% The last argument keeps links to #2.pdf.
\usepackage{xr-hyper}
\newcommand{\maybeexternaldocument}[2]{%
  \IfFileExists{#2.aux}{%
    \externaldocument[#1]{#2}%
  }{%
    \PackageWarning{stack}{Missing #2.aux; external references with prefix #1 unavailable}%
  }%
}
\let\externaldocumentorig\externaldocument
\renewcommand{\externaldocument}[2][]{%
  \IfFileExists{tmp/xref/#2.aux}{%
    \externaldocumentorig[#1]{tmp/xref/#2}[#2.pdf]%
  }{%
  \IfFileExists{#2.aux}{%
    \externaldocumentorig[#1]{#2}%
  }{%
    \PackageWarning{stack}{Missing #2.aux; external references with prefix #1 unavailable}%
  }%
  }%
}

% Package for hypertext links:
\usepackage{hyperref}

% For any local file, say "hello.tex" you want to link to please
% use \externaldocument[hello-]{hello}
\externaldocument[introduction-]{introduction}
\externaldocument[conventions-]{conventions}
\externaldocument[sets-]{sets}
\externaldocument[categories-]{categories}
\externaldocument[topology-]{topology}
\externaldocument[sheaves-]{sheaves}
\externaldocument[sites-]{sites}
\externaldocument[stacks-]{stacks}
\externaldocument[fields-]{fields}
\externaldocument[algebra-]{algebra}
\externaldocument[brauer-]{brauer}
\externaldocument[homology-]{homology}
\externaldocument[derived-]{derived}
\externaldocument[simplicial-]{simplicial}
\externaldocument[more-algebra-]{more-algebra}
\externaldocument[smoothing-]{smoothing}
\externaldocument[modules-]{modules}
\externaldocument[sites-modules-]{sites-modules}
\externaldocument[injectives-]{injectives}
\externaldocument[cohomology-]{cohomology}
\externaldocument[sites-cohomology-]{sites-cohomology}
\externaldocument[dga-]{dga}
\externaldocument[dpa-]{dpa}
\externaldocument[sdga-]{sdga}
\externaldocument[hypercovering-]{hypercovering}
\externaldocument[schemes-]{schemes}
\externaldocument[constructions-]{constructions}
\externaldocument[properties-]{properties}
\externaldocument[morphisms-]{morphisms}
\externaldocument[coherent-]{coherent}
\externaldocument[divisors-]{divisors}
\externaldocument[limits-]{limits}
\externaldocument[varieties-]{varieties}
\externaldocument[topologies-]{topologies}
\externaldocument[descent-]{descent}
\externaldocument[perfect-]{perfect}
\externaldocument[more-morphisms-]{more-morphisms}
\externaldocument[flat-]{flat}
\externaldocument[groupoids-]{groupoids}
\externaldocument[more-groupoids-]{more-groupoids}
\externaldocument[etale-]{etale}
\externaldocument[chow-]{chow}
\externaldocument[intersection-]{intersection}
\externaldocument[pic-]{pic}
\externaldocument[weil-]{weil}
\externaldocument[adequate-]{adequate}
\externaldocument[dualizing-]{dualizing}
\externaldocument[duality-]{duality}
\externaldocument[discriminant-]{discriminant}
\externaldocument[derham-]{derham}
\externaldocument[local-cohomology-]{local-cohomology}
\externaldocument[algebraization-]{algebraization}
\externaldocument[curves-]{curves}
\externaldocument[resolve-]{resolve}
\externaldocument[models-]{models}
\externaldocument[functors-]{functors}
\externaldocument[equiv-]{equiv}
\externaldocument[pione-]{pione}
\externaldocument[etale-cohomology-]{etale-cohomology}
\externaldocument[proetale-]{proetale}
\externaldocument[relative-cycles-]{relative-cycles}
\externaldocument[more-etale-]{more-etale}
\externaldocument[trace-]{trace}
\externaldocument[crystalline-]{crystalline}
\externaldocument[spaces-]{spaces}
\externaldocument[spaces-properties-]{spaces-properties}
\externaldocument[spaces-morphisms-]{spaces-morphisms}
\externaldocument[decent-spaces-]{decent-spaces}
\externaldocument[spaces-cohomology-]{spaces-cohomology}
\externaldocument[spaces-limits-]{spaces-limits}
\externaldocument[spaces-divisors-]{spaces-divisors}
\externaldocument[spaces-over-fields-]{spaces-over-fields}
\externaldocument[spaces-topologies-]{spaces-topologies}
\externaldocument[spaces-descent-]{spaces-descent}
\externaldocument[spaces-perfect-]{spaces-perfect}
\externaldocument[spaces-more-morphisms-]{spaces-more-morphisms}
\externaldocument[spaces-flat-]{spaces-flat}
\externaldocument[spaces-groupoids-]{spaces-groupoids}
\externaldocument[spaces-more-groupoids-]{spaces-more-groupoids}
\externaldocument[bootstrap-]{bootstrap}
\externaldocument[spaces-pushouts-]{spaces-pushouts}
\externaldocument[spaces-chow-]{spaces-chow}
\externaldocument[groupoids-quotients-]{groupoids-quotients}
\externaldocument[spaces-more-cohomology-]{spaces-more-cohomology}
\externaldocument[spaces-simplicial-]{spaces-simplicial}
\externaldocument[spaces-duality-]{spaces-duality}
\externaldocument[formal-spaces-]{formal-spaces}
\externaldocument[restricted-]{restricted}
\externaldocument[spaces-resolve-]{spaces-resolve}
\externaldocument[formal-defos-]{formal-defos}
\externaldocument[defos-]{defos}
\externaldocument[cotangent-]{cotangent}
\externaldocument[examples-defos-]{examples-defos}
\externaldocument[algebraic-]{algebraic}
\externaldocument[examples-stacks-]{examples-stacks}
\externaldocument[stacks-sheaves-]{stacks-sheaves}
\externaldocument[criteria-]{criteria}
\externaldocument[artin-]{artin}
\externaldocument[quot-]{quot}
\externaldocument[stacks-properties-]{stacks-properties}
\externaldocument[stacks-morphisms-]{stacks-morphisms}
\externaldocument[stacks-limits-]{stacks-limits}
\externaldocument[stacks-cohomology-]{stacks-cohomology}
\externaldocument[stacks-perfect-]{stacks-perfect}
\externaldocument[stacks-introduction-]{stacks-introduction}
\externaldocument[stacks-more-morphisms-]{stacks-more-morphisms}
\externaldocument[stacks-geometry-]{stacks-geometry}
\externaldocument[moduli-]{moduli}
\externaldocument[moduli-curves-]{moduli-curves}
\externaldocument[examples-]{examples}
\externaldocument[exercises-]{exercises}
\externaldocument[guide-]{guide}
\externaldocument[desirables-]{desirables}
\externaldocument[coding-]{coding}
\externaldocument[obsolete-]{obsolete}
\externaldocument[fdl-]{fdl}
\externaldocument[index-]{index}

% Heap Project documents
\externaldocument[k3-]{k3}
\externaldocument[torelli-]{torelli}
\externaldocument[linear-algebra-]{linear-algebra}
\externaldocument[groups-]{groups}
\externaldocument[complex-analysis-]{complex-analysis}
\externaldocument[functional-analysis-]{functional-analysis}
\externaldocument[categories-]{categories}
\externaldocument[commutative-algebra-]{commutative-algebra}
\externaldocument[ringed-spaces-]{ringed-spaces}
\externaldocument[homological-algebra-]{homological-algebra}
\externaldocument[lie-algebras-]{lie-algebras}
\externaldocument[representation-theory-]{representation-theory}
\externaldocument[smooth-manifolds-]{smooth-manifolds}
\externaldocument[differential-topology-]{differential-topology}
\externaldocument[algebraic-topology-]{algebraic-topology}
\externaldocument[differential-geometry-]{differential-geometry}
\externaldocument[algebraic-geometry-]{algebraic-geometry}
\externaldocument[symplectic-geometry-]{symplectic-geometry}
\externaldocument[complex-geometry-]{complex-geometry}
\externaldocument[hodge-theory-]{hodge-theory}
\externaldocument[%
variations-of-hodge-structures-]{%
variations-of-hodge-structures}
\externaldocument[%
homological-mirror-symmetry-]{%
homological-mirror-symmetry}
\maybeexternaldocument{deformations-}{deformations}
\externaldocument[vertex-algebras-]{vertex-algebras}
\externaldocument[springer-]{springer}
\externaldocument[infty-categories-]{infty-categories}
\externaldocument[seminars-differential-geometry-]{%
seminars-differential-geometry}
\externaldocument[seminars-complex-geometry-]{%
seminars-complex-geometry}
\externaldocument[seminars-algebraic-geometry-]{%
seminars-algebraic-geometry}
\externaldocument[seminars-analysis-pde-]{%
seminars-analysis-pde}
\externaldocument[surfaces-]{surfaces}
\externaldocument[birational-maps-commutative-algebra-]{birational-maps-commutative-algebra}
\externaldocument[cremona-transformations-]{cremona-transformations}
\externaldocument[derived-categories-of-sheaves-]{derived-categories-of-sheaves}
\externaldocument[geometric-invariant-theory-]{geometric-invariant-theory}
\externaldocument[geometric-stability-conditions-and-group-actions-]{geometric-stability-conditions-and-group-actions}
\externaldocument[moduli-spaces-of-sheaves-]{moduli-spaces-of-sheaves}
\externaldocument[bridgeland-stability-general-theory-]{bridgeland-stability-general-theory}
\externaldocument[hilbert-schemes-]{hilbert-schemes}
\externaldocument[noncommutative-abelian-surfaces-and-kummer-type-hyperkahler-varieties-]{noncommutative-abelian-surfaces-and-kummer-type-hyperkahler-varieties}
\externaldocument[mumford-tate-groups-in-hodge-theory-]{mumford-tate-groups-in-hodge-theory}
\externaldocument[hodge-birational-atoms-2026-]{hodge-birational-atoms-2026}
\externaldocument[xxii-egd-2026-]{%
xxii-egd-2026}
\externaldocument[pde-]{pde}
\externaldocument[probability-]{probability}

\externaldocument[linguistics-]{linguistics}

\externaldocument[%
16th-alga-meeting-2026-]{%
16th-alga-meeting-2026}

\externaldocument[%
iv-jornadas-lefschetz-2026-]{%
iv-jornadas-lefschetz-2026}

\externaldocument[reading-the-masters-]{%
reading-the-masters}

\externaldocument[real-analysis-]{%
real-analysis}

\externaldocument[optimal-transport-]{%
optimal-transport}

\externaldocument[history-]{history}

\externaldocument[seiberg-witten-]{seiberg-witten}

\externaldocument[history-philosophy-biology-]{history-philosophy-biology}

\externaldocument[classical-mechanics-]{classical-mechanics}
\externaldocument[electromagnetism-]{electromagnetism}
\externaldocument[waves-and-optics-]{waves-and-optics}
\externaldocument[wif-]{waves-in-fluids-iv}
\externaldocument[thermal-physics-]{thermal-physics}
\externaldocument[quantum-mechanics-]{quantum-mechanics}
\externaldocument[special-relativity-]{special-relativity}
\externaldocument[physics-]{physics}

\externaldocument[chemistry-]{chemistry}

\externaldocument[molecular-cell-biology-]{molecular-cell-biology}

\externaldocument[microbiology-]{microbiology}

\externaldocument[protistology-]{protistology}

\externaldocument[mycology-]{mycology}

\externaldocument[botany-]{botany}

\externaldocument[zoology-]{zoology}

\externaldocument[animal-histology-]{animal-histology}

\externaldocument[animal-physiology-]{animal-physiology}

\externaldocument[systematics-biogeography-]{systematics-biogeography}

\externaldocument[ecology-natural-resources-]{ecology-natural-resources}

\externaldocument[earth-sciences-]{earth-sciences}

\externaldocument[biostatistics-]{biostatistics}

\externaldocument[mathematical-systems-biology-]{mathematical-systems-biology}

\externaldocument[cimpa-el-salvador-2024-]{cimpa-el-salvador-2024}

\externaldocument[anthropology-]{anthropology}

\externaldocument[yang-mills-theory-]{yang-mills-theory}

\externaldocument[seminars-algebra-]{seminars-algebra}

\externaldocument[seminars-physics-]{seminars-physics}

\externaldocument[seminars-computing-]{seminars-computing}

\externaldocument[oktobermat-2026-]{oktobermat-2026}

\externaldocument[saet-2026-]{saet-2026}

\externaldocument[seminars-symplectic-geometry-]{seminars-symplectic-geometry}

\externaldocument[seminars-dynamics-]{seminars-dynamics}

\externaldocument[oktobermat-2024-]{oktobermat-2024}

\externaldocument[quantum-groups-]{quantum-groups}

\externaldocument[legal-2027-]{legal-2027}

\externaldocument[derived-categories-]{derived-categories}

% Theorem environments.
%
\theoremstyle{plain}
\newtheorem{theorem}[subsection]{Theorem}
\newtheorem{proposition}[subsection]{Proposition}
\newtheorem{lemma}[subsection]{Lemma}
\newtheorem{corollary}[subsection]{Corollary}

\theoremstyle{definition}
\newtheorem{definition}[subsection]{Definition}
\newtheorem{example}[subsection]{Example}
\newtheorem{exercise}[subsection]{Exercise}
\newtheorem{situation}[subsection]{Situation}

\theoremstyle{remark}
\newtheorem{remark}[subsection]{Remark}
\newtheorem{remarks}[subsection]{Remarks}

\numberwithin{equation}{subsection}

% Macros
%
\def\lim{\mathop{\mathrm{lim}}\nolimits}
\def\colim{\mathop{\mathrm{colim}}\nolimits}
\def\Spec{\mathop{\mathrm{Spec}}}
\def\Hom{\mathop{\mathrm{Hom}}\nolimits}
\def\Ext{\mathop{\mathrm{Ext}}\nolimits}
\def\SheafHom{\mathop{\mathcal{H}\!\mathit{om}}\nolimits}
\def\SheafExt{\mathop{\mathcal{E}\!\mathit{xt}}\nolimits}
\def\Sch{\mathit{Sch}}
\def\Mor{\mathop{\mathrm{Mor}}\nolimits}
\def\Ob{\mathop{\mathrm{Ob}}\nolimits}
\def\Sh{\mathop{\mathit{Sh}}\nolimits}
\def\NL{\mathop{N\!L}\nolimits}
\def\CH{\mathop{\mathrm{CH}}\nolimits}
\def\proetale{{pro\text{-}\acute{e}tale}}
\def\etale{{\acute{e}tale}}
\def\QCoh{\mathit{QCoh}}
\def\Ker{\mathop{\mathrm{Ker}}}
\def\Im{\mathop{\mathrm{Im}}}
\def\Coker{\mathop{\mathrm{Coker}}}
\def\Coim{\mathop{\mathrm{Coim}}}

% Boxtimes
%
\DeclareMathSymbol{\boxtimes}{\mathbin}{AMSa}{"02}

%
% Macros for moduli stacks/spaces
%
\def\QCohstack{\mathcal{QC}\!\mathit{oh}}
\def\Cohstack{\mathcal{C}\!\mathit{oh}}
\def\Spacesstack{\mathcal{S}\!\mathit{paces}}
\def\Quotfunctor{\mathrm{Quot}}
\def\Hilbfunctor{\mathrm{Hilb}}
\def\Curvesstack{\mathcal{C}\!\mathit{urves}}
\def\Polarizedstack{\mathcal{P}\!\mathit{olarized}}
\def\Complexesstack{\mathcal{C}\!\mathit{omplexes}}
% \Pic is the operator that assigns to X its picard group, usage \Pic(X)
% \Picardstack_{X/B} denotes the Picard stack of X over B
% \Picardfunctor_{X/B} denotes the Picard functor of X over B
\def\Pic{\mathop{\mathrm{Pic}}\nolimits}
\def\Picardstack{\mathcal{P}\!\mathit{ic}}
\def\Picardfunctor{\mathrm{Pic}}
\def\Deformationcategory{\mathcal{D}\!\mathit{ef}}

\usepackage{graphicx}
\usepackage{url}

\begin{document}

\title{Hilbert schemes}
\author{Daniel González Casanova Azuela}
\maketitle
\phantomsection
\label{section-phantom}

\noindent
This notebook starts with the Hilbert
scheme and its description through families.
It then treats group-invariant subschemes,
generalized Kummer varieties, and seminar
notes on nested Hilbert schemes.
Section \ref{section-intro} uses incidence
correspondences to construct Nakajima's
Heisenberg representation.

\tableofcontents

\section{Hilbert scheme of points}
\label{section-hilbert-scheme}

\noindent
In this section we define
the Hilbert scheme of points
on a smooth projective complex surface $X$.
The essential property of the Hilbert scheme
is that its points (over $\mathbb{C}$)
are 0-dimensional closed subschemes of $X$
of length $n$.

\medskip\noindent
Consider a contravariant functor
$$
\mathcal{H}ilb_X:(\text{Schemes})^{\text{op}}\to (\text{Sets})
$$
which assigns to any scheme $U$ the set of
flat families
of closed subschemes in $X$ parametrized by
$U$.
Any such family is by definition a flat
morphism
$$
\xymatrix@C=.1em{
Z \ar[d]^\pi&\subseteq X \times U\\
U.
}
$$
We define $\mathcal{H}ilb$ to map a morphism
of schemes
$U \to U'$ to the map of sets (of families
over $U$ and $U'$) given
by pullback. That is, we map a flat family
$Z'\to U'$
to $Z'\times_{U'}U \to U$, which remains
flat.

The Hilbert polynomial at a fibre $Z_u$ is
defined by
$P_u(m)=\chi(\mathcal{O}_{Z_u}\otimes
\mathcal{O}_X(m))$.
By flatness of $\pi$ the Hilbert polynomial
is constant along $U$
(\cite[Part III, Theorem 9.9]{har}).
For a fixed polynomial $P$
we define $\mathcal{H}ilb^P_X$ as the
functor which maps $U$
to the set of families over $U$ with Hilbert
polynomial $P$.
A theorem by Grothendieck says that
$\mathcal{H}ilb^P_X$ is representable
by a projective scheme $\text{Hilb}^P_X$.
This means that there is a natural
transformation
$\Hom(-,\text{Hilb}^P_X) \simeq
\mathcal{H}ilb^P_X(-)$;
that is, families over any scheme $U$ are in
bijection with
morphisms $U \to \text{Hilb}^P_X$.

As with any representable functor,
we obtain a distinguished element associated
to the identity morphism
of the object representing the functor.
In our case, this is a
``universal'' family $\mathcal{Z} \to
\text{Hilb}^P_X$
with the property that any other family $Z
\to U$ is obtained as a pullback
of $\mathcal{Z}$.
More exactly, when regarding an arbitrary
family $Z \to U$ as a morphism
$f \in \Hom(U,\text{Hilb}^P_X)$, we see that
$$
\xymatrix{
\text{id}_{\text{Hilb}^P_X}\ar@{|->}[d]
&\Hom(\text{Hilb}^P_X,\text{Hilb}^P_X)\ar[r]^-{\simeq}\ar[d]_{f^*}
&\mathcal{H}ilb^P_X(\text{Hilb}^P_X)\ar[d]^{U\times_{\text{Hilb}_X^P}-}
&\mathcal{Z}\ar@{|->}[d]\\
f&\Hom(U,\text{Hilb}^P_X)\ar[r]^-{\simeq}
 &\mathcal{H}ilb^P_X(U)&
Z=U\times_{\text{Hilb}_X^P}\mathcal{Z}.
}
$$
commutes by naturality.
This shows that studying families of
subschemes of $X$ parametrized by any
scheme $U$ can be done by studying the
universal family
over the Hilbert scheme.

We can formally view subschemes of $X$
as points of $\operatorname{Hilb}_X^P$ as
follows.
Recall that, as with any other scheme,
points $h$ of $\text{Hilb}_X^P$ correspond
to morphisms
$\Spec(\kappa(h)) \to \text{Hilb}^P_X$,
where $\kappa(h)$ is the residue field of
the point
(see Stacks Project
\href{https://stacks.math.columbia.edu/tag/01J9}{01J9}).
Then for every point $h$ of the Hilbert
scheme
we have a unique subscheme
of $X$ defined as the fibre $X_h$
of the corresponding morphism,
$$
\xymatrix{
X_h\ar[r]\ar[d]\ar@{}[dr]|<{\lrcorner}&\mathcal{Z}\ar[d]
& \subset X \times \text{Hilb}^P_X\\
\Spec(\kappa(h))\ar[r]&\text{Hilb}^P_X.
}
$$
As a rather technical observation that will
not be
addressed later, we note that
to make sure that the fibres are
$\mathbb{C}$-schemes,
we must restrict our analysis to
$\mathbb{C}$-points
of the Hilbert scheme (i.e. when the residue
field is $\mathbb{C}$).

\begin{definition}
\label{definition-hilbert-scheme-of-points}
Let $P$ be a constant polynomial given by
$P(m)=n$ for all $m \in \mathbb{Z}$.
The {\it Hilbert scheme of $n$ points in
$X$} is $X^{[n]}:=\text{Hilb}^P_X$.
\end{definition}

\noindent
This condition forces the dimension of the
fibres of any family
over $X$ to be zero. Indeed, the leading term
of the Hilbert polynomial of a projective
variety of degree $d$
and dimension $s$ is $(dm^s)/s!$, which can
only be constant
for $s=0$ (see \cite[Exercise
1.13]{HarrMorr}).

Further, any 0-dimensional subscheme of $X$
is supported on finitely many points of $X$
(this is consequence of $X$ being Noetherian:
we couldn't have infinitely many points
since they
would give an infinite chain of prime
ideals).
Finally, the {\it length} of any such
subscheme, defined by
 $\dim\Gamma(\mathcal{O}_Z)$
must be $\sum_{x \in \text{supp}
Z}\dim\Gamma(\mathcal{O}_{Z_x})=n$.
(see Stacks Project
\href{https://stacks.math.columbia.edu/tag/0AYT}{0AYT}).

\medskip\noindent
It is now easy to see that any subscheme of
$X$ consisting of $n$ distinct
points is itself a point of
$\text{Hilb}_X^P$.
Indeed, if $Z=\{x_1,\ldots,x_n\}\subset X$,
its
structure sheaf is the direct sum of
skyscraper sheaves
$$
\mathcal{O}_Z=\bigoplus_{i=1}^n
(\text{skyscraper sheaf at $x_i$})
$$
so that $\mathcal{O}_Z \otimes
\mathcal{O}_X(m)\cong\mathcal{O}_Z$
and thus the Hilbert polynomial of $Z$ is
constant $n$.
The latter isomorphism can be verified via
the exact sequence associated to the maximal
ideal
$\mathfrak{m}_{x_i}$
of the stalk $\mathcal{O}_{X,x_i}$, namely
$$
\xymatrix{
0\ar[r]
&\mathfrak{m}_{x_i}\ar[r]
&\mathcal{O}_{X,x_i}\ar[r]
&\mathcal{O}_{X,x_i}/\mathfrak{m}_{x_i}\cong\mathbb{C}\ar[r]
&0,
}
$$
which after tensoring by
$\mathcal{O}_{Z,x_i}
\cong\mathbb{C}\cong\mathcal{O}_{X,x_i}/\mathfrak{m}_{x_i}$
kills the first term.
Twisting by $m$ does not alter this
construction.

\medskip\noindent
However, not every point in $X^{[n]}$ is
given
by a set of $n$ distinct points,
nor by a set of points with multiplicites,
as the following construction suggests.
Define the {\it symmetric power} of $X$ as
$$
S^nX=\underbrace{X \times \ldots \times
X}_{n\text{ copies}}/S_n
$$
where the symmetric group $S_n$ acts by
permutating copies. Thus, a point of $S^nX$
is
an $S_n$-orbit on the cartesian product $X
\times \ldots \times X$,
i.e., an unordered $n$-tuple of points of $X$
with allowed repetition.
In algebraic geometry these are called
``algebraic cycles''
and are denoted as formal sums $\sum_{x \in
X}n_x[x]$
where $\sum_{x \in X}n_x=n$, with $n_x \in
\mathbb{N}$,
and $[x]$ are formal symbols.

By our discussion above on the length of
0-dimensional subschemes of $X$,
we have a map
$$
\begin{aligned}
\pi : X^{[n]} &\longrightarrow S^nX \\
 Z&\longmapsto \sum_{x \in \text{supp}Z}\dim
\Gamma(\mathcal{O}_{Z,x})[x]
\end{aligned}
$$
called the {\it Hilbert-Chow morphism}.
In some cases this is an isomorphism, but
not always.

\begin{theorem}[Fogarty, \cite{Fogarty}]
\label{theorem-fogarty}
Suppose $X$ is nonsingular and $\dim X=2$,
then the following holds.
\begin{enumerate}
\item $X^{[n]}$ is nonsingular and has
dimension $2n$.
\item The Hilbert-Chow morphism is a
resolution of singularities.
\end{enumerate}
\end{theorem}




\section{The G-Hilbert scheme}
\label{section-g-hilb}

\noindent
Start with the Hilbert functor of
Section~\ref{section-hilbert-scheme}:
a base $B$ parametrizes finite flat families
$Z\subset X\times B$ of fixed length.
For a finite group $G$ acting on $X$, take
length $|G|$ and require $Z$ to be
$G$-invariant, with $G$ acting trivially on
$B$. There is one extra condition: on each
fibre, the action on its functions must be
isomorphic to the action on functions on a
free $G$-orbit (an orbit of $|G|$ distinct
points). This is the regular representation
condition; invariance and length alone are
not sufficient. The resulting functor is
represented by $G$-Hilb$(X)$. Its component
containing free orbits records their flat
limits, including collisions, and maps to
$X/G$ by sending a subscheme to the orbit
supporting it. This is why a Hilbert scheme
can improve a quotient: it retains limiting
subscheme information that the quotient
forgets. Whether this map resolves the
singularities requires a theorem for the
action in question. See
Bridgeland--King--Reid,
\emph{Mukai implies McKay}, Section 1,
\href{https://arxiv.org/abs/math/9908027}
{arXiv:math/9908027}.

\section{Generalized Kummer varieties}
\label{section-generalized-kummer}

\noindent
Let $A$ be an abelian surface. Recall from
Section
\ref{section-hilbert-scheme}
the Hilbert-Chow morphism
$
\rho:A^{[n+1]}\longrightarrow S^{n+1}A.
$

\noindent
Since $A$ is a group, there is a natural
summation map
$$
\begin{aligned}
\sigma:S^{n+1}A&\longrightarrow A,\\
\sum_i m_i[a_i]&\longmapsto
\sum_i m_i a_i.
\end{aligned}
$$
Here the sum on the right is taken using the
group law of the abelian surface. The map is
well-defined because addition on $A$ is
commutative, so it does not depend on the
ordering of the points.

\begin{definition}
\label{definition-generalized-kummer}
The {\it generalized Kummer variety}
associated to $A$ is
$
K_n(A)
=
(\sigma\circ\rho)^{-1}(0)
\subset A^{[n+1]}.
$
\end{definition}

\begin{theorem}[Beauville]
\label{theorem-generalized-kummer-hk}
Let $A$ be an abelian surface and let
$n\geq1$. Then $K_n(A)$ is a smooth
projective irreducible holomorphic symplectic
variety of dimension $2n$. In particular it
is a hyperk\"ahler variety.
\end{theorem}

\begin{reference}
This is one of the two standard families
constructed in \cite{beauville-kummer}.
\end{reference}

\begin{example}
\label{example-kummer-fourfold}
For $n=2$, consider the zero-sum locus
$
Z=
\{(a_0,a_1,a_2)\in A^3:
a_0+a_1+a_2=0\}.
$
The map
$
A^2\longrightarrow Z,
\qquad
(a,b)\longmapsto(a,b,-a-b)
$
is an isomorphism. The symmetric group
$S_3$ acts on $Z$ by permuting the three
coordinates, and the zero fibre of
$
S^3A\longrightarrow A
$
is
$
Z/S_3\cong A^2/S_3.
$
The Hilbert-Chow morphism restricts to a
resolution
$
K_2(A)\longrightarrow A^2/S_3.
$
Thus $K_2(A)$ is a hyperk\"ahler fourfold
resolving the singular zero-sum quotient.
\end{example}

\section{Irreducibility of the Hilbert scheme
of points and the class of 2-step
ideals}
\label{section-graffeo-two-step-ideals}

\noindent
Michele Graffeo, SISSA.
Algebraic Geometry Seminar, UFF.
September 29, 2025.

\medskip\noindent
These are rough notes from the talk,
relocated from the algebraic geometry
seminar notebook. Their incomplete formulas
and theorem statements still need review.
They are preserved as a reading record,
not as verified statements.

\medskip\noindent
{\bf Abstract.}
Hilbert schemes of points on a
quasi-projective variety X are classical
objects
in algebraic geometry. Roughly speaking, they
parametrise ideals of a polynomial
ring with complex coefficients having finite
colength. Although Hilbert schemes
always have a distinguished component called
the smoothable component, their
geometry is quite pathological and one of the
main open problems around them
concerns their irreducibility. In a joint
work with Giovenzana, Giovenzana,
Lella we introduce the notion of 2-step
ideals. We show that the loci
parametrising these ideals are in general not
contained in the smoothable
component of the Hilbert scheme, thus
providing new examples of
extra-components.  In my seminar I will
discuss our class of ideals and relate
them to the compressed algebras considered by
Iarrobbino in the eighties.
Finally, I will show how to extend our result
to the nested setting.


\medskip\noindent

As usual, you start with a functor and
prove representativity.

\begin{definition}
\label{definition-hilbert-scheme}
Let $X$ be a smooth
quasi-projective variety and a nonegative
integer $d$.
Define
$$
\begin{aligned}
\underline{Hilb}^d: \text{Sch}_\mathbb{C}
&\longrightarrow \mathsf{Set} \\
B &\longmapsto \left\{Z \hookrightarrow B
\times X:
\substack{B\text{-flat} \\ B\text{-finite}\\
\text{length }d}\right\}
\end{aligned}
$$
\end{definition}

\begin{theorem}[Grothendieck]
\label{theorem-grothendieck}
$\underline{Hilb}^d(X)$ is represented by a
quasi-projective
scheme $Hilb^d(X)$
\end{theorem}

Idea: $\mathbb{C}$-points of $Hilb^d(X)$
are in correspondence with $Z
\overset{\subset}{\text{finite}}X$
of length.

A fat point $Z= \text{Spec}(A)$ $(A,
\mathfrak{m})$ is a
local artinian $\mathbb{C}$-algebra of finite
type.

If $X$ smooth can consider $X=\mathbb{A}^n$,
and it's the same to consider
\begin{itemize}
\item $Z \hookrightarrow \mathbb{A}^n$ of
length $d$.
\item $I \subset \mathbb{C}[x_1,\ldots,x_n]$
of
colength $d$.
\item $A= \mathbb{C}[x_1,\ldots,x_n]/I$ with
$\dim_\mathbb{C}=d$.
\end{itemize}

\begin{definition}
\label{definition-nested-hilbert-scheme}
$X$, $0 \leq d_1 \leq \ldots
d_r=\underline{d}$.
$$
Hilb^{\underline{d}}(X)=
\{Z_1\not\hookrightarrow
\ldots
\not\hookrightarrow Z_r  \not\hookrightarrow
X:
\text{len}(Z_i)=d_i\}
$$
\end{definition}

\begin{theorem}[Hartshorne (r=0), Fogarty,
Kalpan (r>1)]
\label{theorem-hk}
$Hilb^d(X)$ is connected.
\end{theorem}

For $r=1$, $Hilb^d(\mathbb{A}^n)$ is smooth
iff
$n \leq 2$, $d \leq 3$.
If $r>1$, $n=2$, it is smooth iff $r=2$ and
$d_1=d_2=-1$.
If $r>1$ and $n>2$, it is smooth iff $r=2$,
$(d_1,d_2)=\{(1,2),(2,3)\}$.

Iror components.
\begin{itemize}
\item $(r=1)$ $n\geq 4$,
$Hilb^d(\mathbb{A}^3)$
is irreducible iff $d\leq 7$ ($\impliedby$
[E-I], $\implies $ [Mozzola].
$n=3$ irredicuble if $d \leq  11$ (8,9 Sivic;
10 Jardim et al; 11, Jelisievv et
al.)
\item $(r=2)$, $n=2$, irreducible
[G-Rosul-Sebastian].
\item $(r \geq 3)$ $n=2$, there exist $d_1
\leq \ldots \leq d_5$
such that
$Hilb^{\underline{d}}(\mathbb{A}^i)$ is
irreducible [S-R].
\end{itemize}

There is an analogy between classical in
dimension 3
and nested in dimension 2.

Schematic structure.
\begin{itemize}
\item [I] $Hilb^{21}(\mathbb{A}^4)$ has at
least a generically
non-reduced component.
\item Problem: what about $n=3$?
\item Theorem. If $Hilb^{d_1,\ldots
d_r}(\mathbb{A}^i)$ is
irreducible $\implies $
$Hilb^{(1,d_1,\ldots,d_r)}(\mathbb{A}^n$
has a generically non reduced component.
\end{itemize}

Other results.
\begin{itemize}
\item …
\end{itemize}
Hilbert scheme function.
For a local artinian $\mathbb{C}$-algebra of
finite type there is an associated graded
$\text{Gr}_\mathfrak{m}(A)
=\bigoplus_{i \geq
0}\mathfrak{m}^1/\mathfrak{m}^{i+1}$
which allows us to define the Hilbert
function
$$
\begin{aligned}
h_A: \mathbb{Z} &\longrightarrow \mathbb{N}
\\
i &\longmapsto \dim_\mathbb{C}
\mathfrak{m}_i/\mathfrak{m}_{i+1}
\end{aligned}
$$
sistability for graded $A$-modules.
The socol is the annihilator of
the ideal which gives the algebra
e.g. $\mathbb{C}[x,y](x^2,xy,y^2)$ then
$\text{Soc}(A)=\left\langle{}x,y^2\right\rangle{}_\mathbb{C}
$.

\medskip\noindent
\begin{definition}
\label{definition-smoothable-component}
The {\it smoothable component} is
 $$
V_{sm}=\overline{\{[(z_i)_{i=1}^n \in
Hilb^{\underline{d}}(\mathbb{A}^n):
Z_r\text{ is reduced}\}}
$$
\end{definition}

\begin{definition}
\label{definition-elementary-componenet}
$V \subset
Hilb^{\underline{d}}(\mathbb{A}^n)$ is
an {\it elementary component} if
$\forall [Z_1,\ldots,Z_r] \in V$
then $Z_r$ is a fat point.
\end{definition}

Instead of looking for all components of the
Hilbert scheme
we look for elementary components. This may
not be some simple;
sometimes there's infinitely many. There
is no method to find them.

\begin{theorem}[Irrabaro]
\label{theorem-i}
Every irreducible component $V \subset
Hilb^d(\mathbb{A}^n)$
is generically étale locally product of
elementary components.
\end{theorem}

\begin{definition}
\label{definition-}
Let $h:\mathbb{Z} \to \mathbb{N}$ be a
function
with finite support, $|h|=\sum_{ i \in
\mathbb{Z}}h(i)=d$. Define
\begin{itemize}
\item $H_n=\{[A] \in Hilb^d(\mathbb{A}^n):
h_A \equiv h\}$.
This is closed by semicontinuity of the
function; but can be expressed
a by a representable functor.
\item The following also has a canonical
schematic structure
$$
\begin{aligned}
\pi_n:H_n  &\longrightarrow
\mathcal{H}_n=\{[A] \in H_n: A\text{ is
graded}\} \\
A &\longmapsto A=\text{Gr}_{\mathfrak{m}}(A)
\end{aligned}
$$
Where the first is $\not\hookrightarrow $ in
the second.
\end{itemize}
\end{definition}

Recall that
$$
T_{[I]}Hilb^d(\mathbb{A}^n)=\Hom_R(I,R/I)
$$
\begin{theorem}[B-B,J,G G G L]
\label{theorem-b-b-j-g-g-g-l}
There is a graded $[I] \in \mathcal{H}_n$.
Then
$$
\begin{aligned}
T_{[I]}\mathcal{H}_n&=\Hom_R(I,R/I)_{=0}\\
T_{[I]}\pi_n([I])&='>0\\
T_{[I]}H_n&=' \geq 0.
\end{aligned}
$$
\end{theorem}

\begin{definition}
\label{definition-compressed}
$(A,\mathfrak{m}_A)$ is {\it compressed} is
for all  $(A',\mathfrak{m}')$
with $e(A)=e(A')$ and $len(A) \geq  len(A')
$.
\end{definition}

\begin{theorem}
\label{theorem-e}
$E_h \subset H_n$ is $H$-compressed
$E_h \neq  \emptyset$
\end{theorem}

\section{Nakajima's representation: overview}
\label{section-intro}
\label{section-abs}

\noindent
Start with the Hilbert scheme $X^{[n]}$ of
Definition
\ref{definition-hilbert-scheme-of-points}
for a smooth projective complex surface $X$.
An incidence correspondence adds or removes
a subscheme concentrated at one point of
$X$. This changes $n$. The operators
therefore act on the direct sum
$$
\bigoplus_{n\geq 0} H_*(X^{[n]}).
$$
Their graded commutators recover the
intersection pairing on $X$. The vacuum
is the empty subscheme in $X^{[0]}$.
See Theorem \ref{theorem-main} and
\cite[Chapter 8]{Nakajima-lectures}.
Section \ref{section-example} computes the
commutator for adding and removing one
point. Section \ref{section-proof} treats
the contribution of modifications at
distinct points.

\section{Heisenberg algebra}
\label{section-heisenberg}

\noindent
The algebra for this construction uses the
graded vector space $V=H_*(X,\mathbb{C})$.
First recall the oscillator relation and
the rank-one character from Vertex Algebras,
Equations \ref{vertex-algebras-%
equation-bracket-basic} and
\ref{vertex-algebras-%
equation-character-heisenberg}.
A negative mode of index $-m$ increases the
number of points by $m$. For rank one the
vacuum representation is the polynomial ring
with one generator for each positive $m$.
\begin{equation}
\label{equation-symmetric-rank-1}
\bigotimes_{m>0}\mathbb{C}[x_m],
\qquad \deg x_m=m .
\end{equation}
If $V$ is purely even, the Fock space has
one copy of $V$ in each positive weight.
\begin{equation}
\label{equation-symmetric}
\bigotimes_{m>0}\operatorname{Sym}(Vt^{-m}).
\end{equation}
This follows from the Poincar\'e--Birkhoff--
Witt theorem since the negative modes
commute. Its character is computed in
Equation \ref{vertex-algebras-%
equation-character-rank-r}.

For a surface with odd homology, give a
homogeneous class $a$ the parity of its
homological degree. The intersection pairing
is graded symmetric. The nonzero modes and
the even central element $K$ span
\begin{equation}
\label{equation-super-heisenberg}
\bigoplus_{m\ne 0} Vt^m\oplus\mathbb{C}K .
\end{equation}
The ordinary bracket for even classes is
\begin{equation}
\label{equation-bracket-affine}
[at^m,bt^n]=m\delta_{m,-n}(a,b)K .
\end{equation}
For homogeneous $a,b$ the graded bracket
obeys the same mode relation
\begin{equation}
\label{equation-super-bracket}
[at^m,bt^n]_{\rm gr}
=m\delta_{m,-n}(a,b)K .
\end{equation}
Thus the negative modes in odd degree give
exterior factors. An even generator of
weight $m$ contributes $1/(1-q^m)$ to the
ordinary character. An odd generator
contributes $1+q^m$. In the surface
representation a class of degree $j$ in
weight $m$ creates a class of degree
$j+2m-2$ from the vacuum. Writing
$b_j=\dim H_j(X,\mathbb{C})$, the resulting
two-variable character is
\begin{equation}
\label{equation-character}
\prod_{m>0}\prod_{j=0}^4
\bigl(1-(-1)^jt^{j+2m-2}q^m\bigr)
^{(-1)^{j+1}b_j}.
\end{equation}

\noindent
It was proved by Göttsche using the Weil
conjectures in \cite{gottsche1990betti}
(and in \cite[Section
6.2]{Nakajima-lectures} by other methods)
that Equation \ref{equation-character}
in fact equals
$$
\sum_{n=0}^\infty q^nP_t(X^{[n]}),
$$
where the {\it Poincaré polynomials}
$P_t(Z)$ are defined by
$P_t(Z)=\sum_{m\geq 0}b_m(Z)t^m$,
$b_m(Z)=\text{rank}H_m(Z)$ for
any variety $Z$.

The representation
of the Heisenberg algebra we discuss in this
work corresponds
precisely to this version of the Heisenberg
algebra.
Göttsche's formula shows that it is a
highest weight representation,
as explained in \cite[Corollary
8.16]{Nakajima-lectures}.






\section{Correspondences}
\label{section-correspondences}

\noindent
In this section we introduce the main tool
for constructing the operators in Nakajima's
representation
of the Heisenberg algebra.

\medskip\noindent
Let $X_1,X_2$ be smooth irreducible varieties
of dimensions $d_1,d_2$.
Let $Z$ be a $d$-dimensional
closed subvariety of $X_1\times X_2$
such that the projection onto the first
factor
is proper.

The {\it correspondence} associated to $Z
\subset X_1\times X_2$
is the map
\begin{equation}
\label{equation-correspondence}
\begin{aligned}
c_{[Z]}:H_i^{BM}(X_2)&\longrightarrow
H_{i+2d-2d_2}^{BM}(X_1)\\
c&\longmapsto p_{1,*}(p_2^*c\cap [Z])
\end{aligned}
\end{equation}

\noindent
where $p_1$ and $p_2$ are the projections in
the following diagram:
$$
\xymatrix{
&Z\ar[ld]_{p_1}\ar[rd]^{p_2} \subset
X_1\times X_2\\
X_1&&X_2.
}
$$
\noindent
That is, the correspondence is an operator
given by pulling back a class by $p_2$,
intersecting with the fundamental class of
$[Z]$,
and pushing forward by $p_1$.
The pullback here is along the projection of
the ambient product. Intersect there with
$[Z]$, then use properness on its support
for the pushforward. See Algebraic Topology,
Section \ref{algebraic-topology-section-bm},
and \cite[Section 8.2]{Nakajima-lectures}.

\begin{remark}
\label{remark-notes}
\begin{enumerate}
\item
\label{item-any-class}
We can define correspondences more generally
replacing the fundamental class of $Z$
by any $\alpha \in H_{2d}^{BM}(Z)$.
In fact, the operators in our representation
shall be of this more general type.

\item If $Z$ is the graph of a proper
morphism
$f:X_2 \to X_1$, then $c_{[Z]}=f_*$.
This shows that correspondences generalize
the pushforward in BM homology.

\item Since $X_1$ and $X_2$ are smooth,
we may use the isomorphism of BM and
relative singular cohomology
to identify a class in the tensor product
with an endomorphism
as we do in general for $V,W$ vector spaces
via the usual isomorphism
$V \otimes W^* \cong \Hom(W,V)$, $v \otimes
\varphi \mapsto
(w \mapsto \varphi(w)v)$.
In this way the fundamental class
$[Z] \in
\bigoplus_{i+j=\dim_\mathbb{R}Z}H_i(X_1)\otimes H_j^{BM}(X_2)$
immediately gives a map
$$
\bigoplus_{n}H_* (X^{[n]}) \to
\bigoplus_{n}H_* (X^{[n\mp i]})
$$
which, in fact, coincides with Equation
\ref{equation-correspondence}
(see \cite[Section 8.3]{Nakajima-lectures}).

\end{enumerate}
\end{remark}

\medskip\noindent
Now we show how to compose correspondences,
as we shall do in the proof of the main
theorem.

Suppose we are given two subvarieties
$Z_{12} \subset X_1 \times X_2$ and $Z_{23}
\subset X_2 \times X_3$
such that the projections to $X_2$ and $X_3$
are proper.
We define the composition of the
corresponding correspondences by
\begin{equation}
\label{equation-composition}
[Z_{12}]\circ[Z_{23}]:=p_{13,*}(p_{12}^*
[Z_{12}]\cap p_{23}^* [Z_{23}])
\end{equation}
according to the diagram
$$
\xymatrix{
&  X_1 \times X_2 \times
X_3\ar[dl]_{p_{12}}\ar[dr]^{p_{23}}\ar[dd]_{p_{13}}\\
 X_1\times X_2& & X_2\times X_3\\
& X_1\times X_3
}
$$
which is well-defined since, in fact,
the projection $p_{13}(p_{12}^{-1}Z_{12}
\cap p_{23}^{-1} Z_{23})\to X_3$
is proper (see \cite[Section
8.3]{Nakajima-lectures}).




\section{Main construction}
\label{section-main-construction}

\noindent
In this section we define the
correspondences used in Nakajima's
representation
of the Heisenberg algebra and formulate
the theorem which guarantees that it is so.

Fix $i>0$ and $n\geq i$. In the product
$$
X^{[n-i]}\times X^{[n]}\times X
$$
take the triples $(J_1,J_2,x)$ satisfying
$$
J_2\subseteq J_1,\qquad
\operatorname{supp}(J_1/J_2)=\{x\}.
$$
Their incidence variety is the $n$th piece
of $P[i]$. Interchanging the first two
factors gives $P[-i]$. Let $\Pi$ project to
$X$ and let $\varpi$ forget the marked
point. For a homogeneous class $\alpha$ on
$X$, the associated correspondence is
\begin{equation}
\label{equation-correspondence-p}
P_\alpha[j]=\varpi_*
  (\Pi^*\alpha\cap[P[j]]),\qquad j=\pm i.
\end{equation}
Here the pullback and intersection take
place in the ambient product. All
projections on the supports are proper
because $X$ is projective. Positive indices
remove points and negative indices add them
as in \cite[Equation
(8.12)]{Nakajima-lectures}.

\medskip\noindent
Before we continue, we formulate and
prove (with the help of other results,
as explained in \cite[Section
8.3]{Nakajima-lectures})
a lemma that will be used later.

\begin{lemma}
\label{lemma-dimension}
The $n$th piece of $P[i]$ has complex
dimension $2n-i+1$.
\end{lemma}

\begin{proof}
Suppose $i>0$.
Fix $\mathcal{J}_1 \in X^{[n-i]}$.
By Theorem \ref{theorem-fogarty} we know that
$\dim X^{[n-i]}=2(n-i)$.
To find the dimension of $P[i]$
we must add to this number the dimension of
the variety
obtained by varying  $\mathcal{J}_2 \in
X^{[n]}$
satisfying
$\operatorname{supp}(\mathcal{J}_1/
\mathcal{J}_2)=\{x\}$
for some $x \in X$.
The generic case is when
$x \in
X\setminus\text{supp}(\mathcal{O}_X/\mathcal{J}_1)$
(this may be accepted under the explanation
that since $\mathcal{O}_X/\mathcal{J}_1$ is
the structure sheaf
of a subvariety, its support has less
dimension than that of $X$;
a formal proof is provided in \cite[Lemma
8.31]{Nakajima-lectures}).
In fact, the set of such $\mathcal{J}_2$ is
$\pi^{-1}(S^i_{(i)}X)$, where $\pi$ is
the Hilbert-Chow morphism and
$S^i_{(i)}(X)=\{i[x] \in S^i(X):x \in X\}$.
This notation is used in other parts of
\cite{Nakajima-lectures};
in our case just notice that $S^i_{(i)}(X)$
is the subset of the symmetric product
$S^i(X)$
consisting of cycles of a single point with
weight $i$.
Finally, $\pi^{-1}(S^i_{(i)}(X))$ has
complex dimension $i+1$
as a subvariety of the Hilbert scheme.
\end{proof}

\medskip\noindent
For homogeneous classes $\alpha,\beta$ put
$\epsilon=(-1)^{|\alpha||\beta|}$. For
operators $A,B$ of those parities write
$[A,B]_{\rm gr}=AB-\epsilon BA$.

\begin{theorem}[Nakajima, Grojnowski]
\label{theorem-main}
\begin{reference}
\cite[Theorem 8.13]{Nakajima-lectures}
\end{reference}
On the direct sum of the homology groups of
the Hilbert schemes of a smooth projective
complex surface, the incidence operators
satisfy
\begin{equation}
\label{equation-main}
[P_\alpha[i],P_\beta[j]]_{\rm gr}
=(-1)^{i-1}i\delta_{i,-j}
(\alpha,\beta)\operatorname{Id}
\end{equation}
for all nonzero indices $i,j$. Here
$(\alpha,\beta)$ is the degree of the
intersection pairing, extended by zero in
noncomplementary degrees.
\end{theorem}

The bracket is an anticommutator when both
classes are odd. For even classes it is the
ordinary commutator. Multiplying the
creation operator of index $-m$ by
$(-1)^{m-1}$ for $m>0$ puts this relation
in the oscillator normalization of Section
\ref{section-heisenberg}.

\section{Worked example}
\label{section-example}

\noindent
Let $Z$ consist of $n$ distinct points of
$X$, and choose $x$ outside $Z$. Write
$C=P_{[x]}[-1]$ for creation at $x$ and
$D=P_{[X]}[1]$ for removal at any point.
Then $C[Z]=[Z\cup\{x\}]$. Applying $D$
next gives one term for each point of
$Z\cup\{x\}$ that can be removed. Removing
$x$ gives back $[Z]$. Removing $y\in Z$
gives the same result as first removing $y$
and then adding $x$. Thus
$$
\begin{aligned}
DC[Z]&=[Z]+
\sum_{y\in Z}[(Z\setminus\{y\})\cup\{x\}],\\
CD[Z]&=
\sum_{y\in Z}[(Z\setminus\{y\})\cup\{x\}].
\end{aligned}
$$
Subtracting gives $[D,C][Z]=[Z]$.

The two orders of adding and removing a
point are shown in \cite[Example
8.17]{Nakajima-lectures}.
$$
\includegraphics[width=.85\linewidth]{%
slides/talks/nakajima-heisenberg/fig1.png}
$$



\section{Proof of one case of the main
theorem}
\label{section-proof}

\noindent
The proof of Theorem \ref{theorem-main} is
split
in two different cases:
\begin{enumerate}
\item $i,j>0$ or $i,j<0$, and
\label{item-same}
\item $i>0$, $j<0$ or $i<0$, $j>0$.
\label{item-different}
\end{enumerate}

\noindent
In this work we explain the proof of case
(\ref{item-same}).

\medskip\noindent
Suppose $i,j>0$ (the case $i,j<0$ is
analogous).
Recalling our general definition for
composition of correspondences
(Equation \ref{equation-composition})
we see that
the composition $P_\alpha[i]P_\beta[j]$ is
given by
\begin{equation}
\label{equation-composition-pij}
P_\alpha[i]P_\beta[j]
=p_{13,*}(p_{12}^*[P[i]]\cap\Pi_1^*\alpha\cap p_{23}^*[P[j]]\cap\Pi_2^*\beta).
\end{equation}

Consider the set
\begin{equation}
\label{equation-cycle-pij}
\begin{aligned}
P[i,j]&:=p_{12}^{-1}(P[i])\cap
p_{23}^{-1}(P[j])
\cap(X^{[n-i-j]}\times X^{[n-j]}\times
X^{[n]})\\
&=\left\{(\mathcal{J}_1,\mathcal{J}_2,\mathcal{J}_3):
\mathcal{J}_1\supset\mathcal{J}_2\supset\mathcal{J}_3,
\substack{
\text{supp}(\mathcal{J}_1/\mathcal{J}_2)=\{x\}
\text{ for some }x\in X\\
\text{supp}(\mathcal{J}_2/\mathcal{J}_3)=\{y\}
\text{ for some }y \in X }\right\}
\end{aligned}
\end{equation}

\noindent
which corresponds to the diagram
$$
\xymatrix{
&  X^{[n-i-j]} \times X^{[n-j]} \times
X^{[n]}
\ar[dl]_{p_{12}}\ar[dr]^{p_{23}}\ar[dd]_{p_{13}}\\
P[j]\subset X^{[n-j]}\times X^{[n]}& &
X^{[n-i-j]}\times X^{[n-j]}\supset P[i]\\
& X^{[n-i-j]}\times X^{[n]}.
}
$$
By definition, the correspondence $P[j]$
maps the homology of $X^{[n]}$ to the
homology of $X^{[n-j]}$. In order to compose
this with $P[i]$
we must regard $P[i]$ as a map from the
homology of $X^{[n-j]}$
to the homology of $X^{[n-j-i]}$, which is
perfectly
valid since our correspondences are defined
for all $n$.
Observe that the middle factor $X^{[n-j]}$
in the triple product will become
$X^{[n-i]}$ when
we compute the composition in the opposite
order.

\medskip\noindent
To recover the correspondence as in
Equation \ref{equation-composition-pij}
from the set $P[i,j]$ it's natural to
consider its
fundamental class.
We will do this in the following special way.
Observe that $P[i,j]$ may be decomposed as a
disjoint union
of the set $\{x \neq y\}$
of ``generic'' cases in which $x$ and $y$ in
Equation \ref{equation-cycle-pij}
are different, and the set $\{x=y\}$ in
which they coincide.
Let
$\overline{P}[i,j]:=\overline{P[i,j]\cap\{x\neq y\}}$.
Then
$$
p_{12}^* [P[i]]\cap p_{23}^*
[P[j]]=[\overline{P}[i,j]]+\iota_*R
$$
where $R \in H_*^{BM}(P[i,j]\cap \{x=y\})$
corresponds
to the $\{x=y\}$ part of $P[i,j]$
and $\iota$ is the inclusion.

\medskip\noindent
We can carry out the same construction
for the composition $P_\beta[j]P_\alpha[i]$,
obtaining
\begin{equation}
\label{equation-composition-pji}
P_\alpha[j]P_\beta[i]
=p_{13,*}(p_{12}^*[P[j]]\cap\Pi_2^*\alpha\cap p_{23}^*[P[i]]\cap\Pi_1^*\beta).
\end{equation}

\noindent
and
\begin{equation}
\label{equation-pji-r}
p_{12}^* [P[j]]\cap p_{23}^*
[P[i]]=[\overline{P}[j,i]]+\iota'_*R'.
\end{equation}

\noindent
Observe that Equations
\ref{equation-composition-pij} and
\ref{equation-composition-pji} differ in the
classes
$\Pi_1^*\alpha$ against $\Pi_2^*\alpha$, and
correspondingly with $\beta$.
The following lemma is all we need to obtain
the result:

\begin{lemma}
\label{lemma-pij}
\begin{enumerate}
\item There exists an isomorphism
\begin{equation}
\label{equation-pij-isomorphism}
P[i,j]\cap\{x \neq
y\}\xrightarrow{\cong}P[j,i]\cap \{x \neq
y\},
\end{equation}

\noindent
which exchanges $\Pi_1$ and $\Pi_2$.

\item We have
\begin{equation}
\label{equation-r}
p_{13,*}(\Pi_1^*\alpha\cap\Pi_2^*\beta\cap\iota_*R)=0,\qquad
p_{13,*}(\Pi_1\beta\cap\Pi_2^*\alpha\cap\iota_*R')=0.
\end{equation}
\end{enumerate}
\end{lemma}

\noindent
Indeed, this lemma is all we need since
$$
\begin{aligned}
P_\alpha[i]P_\beta[j]
&=(-1)^{\text{deg}\alpha}p_{13,*}((\overline{P}[i,j]+\iota_*R)\cap\Pi_1^*\alpha\cap\Pi_2^*\beta)\\
&=(-1)^{\text{deg}\alpha\text{deg}\beta}
p_{13,*}((\overline{P}[j,i]+\iota_*R')\cap\Pi_1^*\beta\cap\Pi_2^*\alpha)\\
&=P_\beta[j]P_\alpha[i].
\end{aligned}
$$

\begin{proof}[Proof of Lemma \ref{lemma-pij}]
\begin{enumerate}
\item In this proof I assume that
``isomorphism'' means bijection, since
$P[i,j]$ is considered as a set-theoretical
intersection.

First notice that when we say that the
isomorphism, say $\varphi$,
``exchanges $\Pi_1$ and $\Pi_2$''
we mean that $\Pi_1|_{X \times X}=\Pi_2|_{X
\times X}\circ\varphi$
according to the diagram
$$
\xymatrix{
& & X\\
P[i,j] \cap \{x\neq y\}\subset
\ar[ddd]_\varphi&
\save{
X^{[n-i-j]}\times X^{[n-j]} \times X^{[n]}
\times X \times X
}
\ar[ur]^{\Pi_1}\ar[dr]_{\Pi_2}
\restore
\\
& & X\\
& & X\\
P[j,i]\cap\{x\neq y\}\subset&
X^{[n-i-j]}\times X^{[n-i]} \times X^{[n]}
\times X \times X
\ar[ur]^{\Pi_1}\ar[dr]_{\Pi_2}\\
& & X.
}
$$

Now, the point of the isomorphism $\varphi$
is that it must
map an ideal in $X^{[n-j]}$ to an ideal in
$X^{[n-i]}$.
More precisely, we map
$(\mathcal{J}_1,\mathcal{J}_2,\mathcal{J}_3)\in P[i,j] \cap \{x \neq  y\}$
to
$(\mathcal{J}_1,\mathcal{J}_2',\mathcal{J}_3) \in P[j,i]\cap \{x \neq y\}$.
The defining relation is
\begin{equation}
\label{equation-isomorphism}
\mathcal{J}_1/\mathcal{J}_2'=\mathcal{J}_2/\mathcal{J}_3,\qquad
\mathcal{J}_2'/\mathcal{J}_3=\mathcal{J}_1/\mathcal{J}_2.
\end{equation}

\noindent
The diagram from \cite[Section
8.4]{Nakajima-lectures} shows the two
intermediate ideals.
$$
\includegraphics[width=.8\linewidth]{%
slides/talks/nakajima-heisenberg/fig2.png}
$$

The setting is the following. We have the
ideals with contentions
$\mathcal{J}_1\supset
\mathcal{J}_2,\mathcal{J}_2'\supset
\mathcal{J}_3$.
Passing to the supports, the inclusions are
reversed, i.e
the support of $\mathcal{J}_3$ contains the
rest,
and the support of $\mathcal{J}_1$ is
contained in the rest.

To prove that this function is well-defined
we need to show that for every
$\mathcal{J}_2$ there is a unique
$\mathcal{J}_2'$ satisfying Equation
\ref{equation-isomorphism}.
Notice that by definition,
$\text{supp}\mathcal{J}_2/\mathcal{J}_3=\{x\}$
for some $x \in X$ so that, as a scheme,
$\{x\}$ has multiplicity $j$.
We define
$\mathcal{J}_2'=\mathcal{J}_1\cup \{x\}$
where $x$ has multiplicity $j$.

\medskip\noindent
(Caveat: it looks to me that the choice of
$\mathcal{J}_2'$ is not unique.
Indeed, the space of schemes of length $j$
supported
in the point $x$ was discussed in the proof
of Lemma \ref{lemma-dimension},
where we used the fact that it has dimension
$j-1$.)

\medskip\noindent
Defining $\mathcal{J}'_2$ this way implies
that
$$
\mathcal{J}_1/\mathcal{J}_2'=\mathcal{J}_1/(\mathcal{J}_1\cup \{x\})
=\{x\}=\mathcal{J}_2/\mathcal{J}_3.
$$
Now, we also must have
$\mathcal{J}_1/\mathcal{J}_2=\{y\}$
for some $y \in X$ with multiplicity $i$.
This implies that
$\mathcal{J}_1/\mathcal{J}_3=\{x\}\cup
\{y\}$
with their multiplicities.
Then
$$
\mathcal{J}_2'/\mathcal{J}_3=(\mathcal{J}_1\cup \{x\})/\mathcal{J}_3=\{y\}
=\mathcal{J}_1/\mathcal{J}_2
$$
as required.
An inverse function
defined similarly shows this is a bijection.

\item Consider the commutative diagram
$$
\xymatrix{
P[i,j]\cap
\{x=y\}\ar[r]^-{\Pi'}\ar[d]_{\iota}&X\ar[d]^{\Delta}\\
P[i,j]\ar[r]_{\Pi_1 \times \Pi_2}&X \times X
}
$$
where $\Delta$ is the diagonal embedding
and $\Pi'$ is any of the projections $\Pi_1$
or $\Pi_2$
(which coincide on  $P[i,j]\cap \{x=y\}$).
Then
\begin{equation}
\label{equation-lemma2-1}
\begin{aligned}
\Pi_1^*\alpha \cap \Pi_2^*\beta \cap\iota_*R
&=\iota_*(\iota^*(\Pi_1^*\alpha\cap\Pi_2^*\beta)\cap R)\\
&=\iota_*(\iota^*(\Pi_1\times\Pi_2)^*(\alpha\times\beta)\cap R)\\
&=\iota_*({\Pi'}^*\Delta^*(\alpha\times\beta)\cap R)\\
&=\iota_*({\Pi'}^*(\alpha\cap\beta)\cap R),
\end{aligned}
\end{equation}

\noindent
where the first equality is due
to the projection formula
(see Stacks Project
\href{https://stacks.math.columbia.edu/tag/01E8}{01E8})
and we stick with Nakajima's notation
$\alpha \times \beta$
for the tensor product $\alpha \otimes
\beta$.
We also have the commutative diagram
$$
\xymatrix{
p_{13}\circ \iota:P[i,j]\cap\{x=y\}
\ar[r]^-{\Pi'}\ar[d]_{p_{13 \circ
\iota}}&X\ar@{=}[d]\\
P[i+j]=\{(\mathcal{J}_1,\mathcal{J}_3)|\mathcal{J}_1 \supset \mathcal{J}_3,
\text{supp}(\mathcal{J}_1/\mathcal{J}_3)=\{x\}\}\ar[r]_-{\Pi''}&X,
}
$$
where $\Pi''$ is the usual projection to $X$
from the
set $P[i+j]$.

By Equation \ref{equation-lemma2-1}
and the commutativity of the second diagram,
we can express the class of interest as
intersection
of two classes: one pulled back directly
from $X$,
and the other one corresponding to $R$.
That is, we have
$$
p_{13,*}\iota_*({\Pi'}^* (\alpha\cap
\beta)\cap R)
={\Pi''}^* (\alpha \cap \beta)\cap
p_{13,*}\iota_*(R).
$$
Now, by Lemma \ref{lemma-dimension},
we know that
$$
\dim_{\mathbb{C}}P[i+j]=2n-i-j+1.
$$
On the other hand, the ``expected dimension''
is $2n-i-j+2$. This obliges $\iota_*R$ to
vanish.

\medskip\noindent
(Caveat: I understand that the ``expected
dimension''
means in \cite{Nakajima-lectures} the
dimension of
$P[i,j]\cap \{x\neq y\}$, which is claimed
to be
$2n-i-j+2$.
It is however not clear to me why this would
imply
that $p_{13,*}\iota_*R$ must vanish.)

\end{enumerate}
\end{proof}

\bibliography{my}
\bibliographystyle{amsalpha}

\end{document}
